\documentclass[11pt,a4paper]{article}

% =================================================
% Paquetes
% =================================================

\usepackage[utf8]{inputenc}
\usepackage[T1]{fontenc}
\usepackage[spanish]{babel}
\usepackage{geometry}
\usepackage{enumitem}
\usepackage{hyperref}
\usepackage{xcolor}
\usepackage{listings}
\usepackage{booktabs}
\usepackage{array}
\usepackage{tabularx}
\usepackage{parskip}

\geometry{
  top=2.5cm,
  bottom=2.5cm,
  left=2.7cm,
  right=2.7cm
}

\hypersetup{
  colorlinks=true,
  linkcolor=black,
  urlcolor=blue
}

\setlist[itemize]{
  noitemsep,
  topsep=4pt
}

\setlist[enumerate]{
  noitemsep,
  topsep=4pt
}

\lstset{
  basicstyle=\ttfamily\small,
  frame=single,
  breaklines=true,
  columns=fullflexible,
  showstringspaces=false
}

\lstset{
  literate=
  {á}{{\'a}}1 {é}{{\'e}}1 {í}{{\'i}}1 {ó}{{\'o}}1 {ú}{{\'u}}1
  {Á}{{\'A}}1 {É}{{\'E}}1 {Í}{{\'I}}1 {Ó}{{\'O}}1 {Ú}{{\'U}}1
  {ñ}{{\~n}}1 {Ñ}{{\~N}}1 {¿}{{?`}}1 {¡}{{!`}}1
}

% =================================================
% Documento
% =================================================

\title{
  \textbf{Guía de Tecnologías, Arquitectura y Despliegue}\\
  \large Taller de Integración II y Taller de Integración IV
}

\author{}
\date{Segundo semestre 2026}

\begin{document}

\maketitle

% =================================================
\section{Objetivo}
% =================================================

Esta guía reúne el stack, la arquitectura, la estructura del monorepo y los
procedimientos básicos de desarrollo y despliegue de Taller de Integración II
(TI2) y Taller de Integración IV (TI4).

La guía sirve para:

\begin{itemize}
  \item establecer un stack tecnológico común;
  \item explicar la responsabilidad de cada tecnología;
  \item justificar las principales decisiones tecnológicas;
  \item definir la estructura inicial del monorepo;
  \item explicar cómo ejecutar los proyectos localmente;
  \item definir cómo se comunican las aplicaciones con el backend;
  \item establecer cómo se gestionan variables de entorno y secretos;
  \item explicar cómo desplegar las diferentes partes del sistema;
  \item establecer claramente las responsabilidades técnicas de TI2 y TI4;
  \item reducir decisiones técnicas innecesarias durante el desarrollo.
\end{itemize}

La guía complementa el protocolo de trabajo con Git, GitHub y Linear.

% =================================================
\section{Principio general de arquitectura}
% =================================================

El producto vive en un único \textbf{monorepo GitHub}. El monorepo contiene las
aplicaciones Web y Mobile, el backend compartido y los paquetes internos que se
incorporen durante el desarrollo.

La estructura de carpetas no basta para describir la arquitectura. El backend y
ambas aplicaciones cliente siguen una \textbf{arquitectura por capas}.

\begin{lstlisting}
Proyecto
|
+-- apps/
|   +-- web/                 Aplicación Web de TI2
|   +-- mobile/              Aplicación Mobile de TI4
|
+-- convex/                  Backend, base de datos y autenticación
+-- packages/                Código compartido cuando corresponda
+-- .github/workflows/       Integración continua
+-- package.json             Configuración de workspaces
+-- bun.lock                 Único archivo de bloqueo
\end{lstlisting}

Web y Mobile comparten:

\begin{itemize}
  \item el mismo backend;
  \item la misma base de datos;
  \item el mismo sistema de autenticación;
  \item las mismas reglas de negocio.
\end{itemize}

La arquitectura lógica es:

\begin{lstlisting}
                 PRESENTACIÓN
       Web y Mobile + funciones publicas
                         |
                         v
                    APLICACIÓN
                 Casos de uso
                         |
                         v
                     DOMINIO
             Reglas y políticas del negocio
                         ^
                         |
                 INFRAESTRUCTURA
        Convex, Better Auth y servicios externos
\end{lstlisting}

Una solicitud sigue la dirección \textbf{Presentación -> Aplicación ->
Dominio}. Infraestructura implementa las interfaces que necesitan Aplicación y
Dominio, por lo que puede depender de esas capas. Las reglas del negocio no
dependen de Convex, React, Expo ni Better Auth.

\subsection{Responsabilidad de cada capa}

\begin{description}
  \item[Presentación:] recibe interacciones, valida el formato de entrada,
    convierte la solicitud en una llamada al caso de uso y presenta el
    resultado. Incluye pantallas, componentes, routers y funciones públicas de
    Convex.
  \item[Aplicación:] implementa casos de uso y coordina el flujo de una
    operación. Decide qué servicios necesita, controla la autorización de la
    operación y define los puertos que utilizará para acceder a recursos
    externos.
  \item[Dominio:] contiene entidades, objetos de valor, políticas y reglas de
    negocio. Esta capa debe poder probarse sin levantar React, Expo, Convex ni
    una base de datos.
  \item[Infraestructura:] contiene adaptadores concretos para persistencia,
    autenticación, servicios externos y configuración de ejecución. Aquí se
    ubican el esquema de Convex, los repositorios, Better Auth y el código
    generado.
\end{description}

\subsection{Reglas de dependencia}

Estas reglas rigen las dependencias:

\begin{itemize}
  \item una capa externa puede utilizar una capa interna, pero no al revés;
  \item Presentación no accede directamente a la base de datos;
  \item Aplicación no contiene componentes React ni consultas directas a
    Convex;
  \item Dominio no importa librerías de interfaz, infraestructura, variables
    de entorno ni tipos generados;
  \item Infraestructura implementa puertos definidos por Aplicación o Dominio;
  \item Web y Mobile no duplican reglas de negocio que ya pertenecen al
    backend;
  \item los paquetes compartidos contienen contratos, tipos o utilidades
    reutilizables, pero no se convertirán en un depósito de lógica sin dueño.
\end{itemize}

TI2 es propietario técnico del backend compartido y de sus capas de Aplicación,
Dominio e Infraestructura. TI4 consume esos servicios y mantiene la
Presentación Mobile desde la misma base de código.

% =================================================
\section{Stack tecnológico}
% =================================================

El stack es:

\begin{table}[h]
  \centering
  \begin{tabularx}{\textwidth}{>{\bfseries}l X}
    \toprule
    Componente & Tecnología \\
    \midrule
    Modelo arquitectónico & Arquitectura por capas \\
    Lenguaje principal & TypeScript \\
    Frontend Web & React \\
    Build Web & Vite \\
    Routing Web & TanStack Router \\
    Aplicación móvil & React Native + Expo \\
    Routing móvil & Expo Router \\
    Backend & Convex \\
    Base de Datos & Convex \\
    Autenticación & Better Auth \\
    Build móvil & EAS Build \\
    Hosting Web & Vercel \\
    Hosting Backend & Convex Cloud \\
    Control de versiones & Git + GitHub \\
    Gestión del trabajo & Linear \\
    Integración continua & GitHub Actions \\
    Organización del código & Monorepo con Bun workspaces \\
    \bottomrule
  \end{tabularx}
\end{table}

Ejecuta los comandos con \texttt{bun} desde la raíz del monorepo, salvo cuando
la guía indique explícitamente una aplicación.

Mantén un único gestor de paquetes y un único archivo de bloqueo en la raíz del
monorepo. Las aplicaciones y los paquetes internos usan los workspaces
definidos en el \texttt{package.json} raíz.

\subsection{Instalación de Bun}

Antes de ejecutar la guía, instala Bun y verifica que esté disponible en la
terminal.

En macOS y Linux, ejecuta:

\begin{lstlisting}
curl -fsSL https://bun.sh/install | bash
\end{lstlisting}

En Windows, ejecuta PowerShell:

\begin{lstlisting}
powershell -c "irm bun.sh/install.ps1|iex"
\end{lstlisting}

La mayoría de TI2 usa Windows, por lo que esa es la instalación de referencia
para el equipo. Después de instalar Bun, abre una nueva terminal y comprueba la
instalación con:

\begin{lstlisting}
bun --version
\end{lstlisting}

En el monorepo no deben coexistir:

\begin{lstlisting}
            bun.lock
            package-lock.json
            pnpm-lock.yaml
            yarn.lock
\end{lstlisting}

% =================================================
\section{TypeScript}
% =================================================

\subsection{Responsabilidad}

TypeScript es el lenguaje principal de ambos equipos. Se usa en:

\begin{itemize}
  \item frontend web;
  \item aplicación móvil;
  \item funciones del backend;
  \item definición del modelo de datos;
  \item autenticación;
  \item contratos entre clientes y backend.
\end{itemize}

\subsection{Justificación}

Usar TypeScript en todo el proyecto permite compartir tipos entre equipos y
detectar muchos errores antes de ejecutar la aplicación.

Además, Convex permite obtener tipos de extremo a extremo entre sus funciones
de backend y los clientes TypeScript.

Escribe el código nuevo en TypeScript.

% =================================================
\section{React}
% =================================================

React construye la interfaz de usuario de la aplicación web de TI2 y pertenece
principalmente a la capa de Presentación. React maneja:

\begin{itemize}
  \item componentes visuales;
  \item formularios;
  \item composición de pantallas;
  \item interacción con el usuario;
  \item presentación de los datos obtenidos desde los casos de uso;
  \item captura de eventos y transformación de acciones del usuario en
    comandos de Aplicación.
\end{itemize}

React no accede directamente a la base de datos ni contiene reglas de negocio
que deban cumplirse de la misma forma en Web y Mobile.

El flujo esperado es:

\begin{lstlisting}
Componente React
      |
      v
Caso de uso del cliente
      |
      v
Cliente Convex
      |
      v
Función pública de Convex
      |
      v
Capas de Aplicación y Dominio
\end{lstlisting}

% =================================================
\section{Vite}
% =================================================

Vite ejecuta el entorno de desarrollo y construye la aplicación web. Se ocupa
de:

\begin{itemize}
  \item ejecutar el servidor de desarrollo;
  \item procesar TypeScript y JSX;
  \item administrar módulos durante desarrollo;
  \item generar el build de producción.
\end{itemize}

\subsection{Justificación}

La aplicación web funciona principalmente como cliente de Convex. Vite se
ocupa del desarrollo y la construcción, no de la lógica de negocio. Por eso no
se incorpora un segundo framework full-stack con otra capa de servidor.

La distribución del cliente Web queda así:

\begin{lstlisting}
Presentación
 |
 +-- React
 +-- TanStack Router
 |
 +-- Aplicación del cliente
       |
       +-- Infraestructura del cliente
             |
             +-- Convex Client
             +-- Better Auth Client
\end{lstlisting}

La lógica de servidor permanecerá en Convex y seguirá la separación entre
Presentación, Aplicación, Dominio e Infraestructura definida anteriormente.

% =================================================
\section{TanStack Router}
% =================================================

TanStack Router administra las rutas de la aplicación web. Se ocupa de:

\begin{itemize}
  \item definición de páginas;
  \item navegación;
  \item parámetros de URL;
  \item layouts;
  \item rutas anidadas;
  \item navegación type-safe.
\end{itemize}

Las rutas se definen mediante archivos.

Por defecto las rutas se almacenarán en:

\begin{lstlisting}
src/routes/
\end{lstlisting}

TanStack Router genera automáticamente:

\begin{lstlisting}
src/routeTree.gen.ts
\end{lstlisting}

No edites este archivo manualmente.

\subsection{Justificación}

Se usa TanStack Router en lugar de TanStack Start porque Convex ya define el
servidor del proyecto. TanStack Start incorporaría funcionalidades de servidor
que la arquitectura no necesita.

% =================================================
\section{Convex}
% =================================================

Convex es la plataforma central del backend. Aporta:

\begin{itemize}
  \item backend;
  \item base de datos;
  \item plataforma de ejecución de funciones;
  \item mecanismo de sincronización de datos con los clientes.
\end{itemize}

TI2 desarrolla y mantiene el backend.

Convex no define por sí solo la arquitectura. Sus funciones públicas actúan
como adaptadores de Presentación: reciben la solicitud, validan su forma y
delegan en un caso de uso. No concentran consultas, autorización y reglas de
negocio en una sola función extensa.

\subsection{Tipos de funciones}

La API pública del backend usa:

\begin{description}
  \item[Queries:] operaciones de lectura.
  \item[Mutations:] operaciones que modifican datos.
  \item[Actions:] operaciones que requieren interacción con servicios externos
    o comportamientos que no corresponden a una query o mutation.
\end{description}

Queries, mutations y actions forman el borde de entrada del backend. El caso de
uso que invocan pertenece a Aplicación; las reglas que debe respetar pertenecen
al Dominio.

La API se distribuye así:

\begin{lstlisting}
Cliente
  |
  +-- Query -------> Leer
  |
  +-- Mutation ----> Modificar
  |
  +-- Action ------> Servicio externo / proceso especial
\end{lstlisting}

La lógica de negocio no debe quedar directamente en queries, mutations o
actions. Mantén delgadas estas funciones para poder probar y reutilizar el caso
de uso sin depender del runtime de Convex.

% =================================================
\section{Base de datos Convex}
% =================================================

Convex administra el modelo de datos, que pertenece a Infraestructura.

El esquema principal se encontrará en:

\begin{lstlisting}
convex/schema.ts
\end{lstlisting}

El esquema debe estar bajo control de versiones.

El esquema describe cómo se persisten los datos, pero no reemplaza las
entidades ni las reglas del Dominio. Cuando una regla de negocio necesita leer o
guardar información, Aplicación usa un puerto e Infraestructura lo implementa
con Convex.

Para cada cambio relevante del modelo de datos:

\begin{itemize}
  \item realizarse mediante una branch;
  \item pasar por Pull Request;
  \item considerar el impacto sobre la aplicación web;
  \item considerar el impacto sobre la aplicación móvil;
  \item quedar asociados al issue correspondiente en Linear.
\end{itemize}

No dejes cambios manuales indispensables que solo conozca una persona.

% =================================================
\section{Better Auth}
% =================================================

Better Auth administra la autenticación de usuarios y pertenece a
Infraestructura del backend.

La instancia principal de Better Auth corre sobre Convex.

Sus responsabilidades incluyen:

\begin{itemize}
  \item registro de usuarios;
  \item inicio de sesión;
  \item cierre de sesión;
  \item manejo de sesiones;
  \item mecanismos adicionales de autenticación habilitados por el proyecto.
\end{itemize}

El adaptador de Better Auth entrega la identidad autenticada al borde de
entrada. Aplicación y Dominio usan esa identidad para evaluar autorización y
políticas del negocio, sin importar directamente módulos de Better Auth.

\subsection{Autenticación y autorización}

Distingue:

\begin{description}
  \item[Autenticación:] determina quién es el usuario.
  \item[Autorización:] determina qué puede hacer el usuario.
\end{description}

En consecuencia:

\begin{lstlisting}
Infraestructura
    |
    +--> Better Auth entrega la identidad
    |
    v
Presentación
    |
    v
Aplicación
    |
¿Puede ejecutar esta operación?
    |
    v
Dominio evalúa la política de autorización
\end{lstlisting}

Implementa las reglas de autorización del negocio en Aplicación o Dominio,
según corresponda. No confíes solo en ocultar una pantalla del frontend.

% =================================================
\section{Expo y React Native}
% =================================================

TI4 desarrolla la aplicación móvil con React Native mediante Expo. React Native
pertenece a Presentación Mobile y Expo proporciona el entorno de ejecución.

Expo proporciona:

\begin{itemize}
  \item entorno de desarrollo móvil;
  \item herramientas de ejecución;
  \item configuración de Android e iOS;
  \item integración con APIs nativas;
  \item proceso de build mediante EAS.
\end{itemize}

\subsection{Justificación}

Expo reduce la configuración nativa que el equipo administra directamente y
mantiene la mayor parte del desarrollo en TypeScript y React. Así, ambos
equipos comparten tecnologías y conceptos.

\begin{center}
  \begin{tabular}{ll}
    \toprule
    TI2 Web & TI4 Mobile \\
    \midrule
    React & React Native \\
    TanStack Router & Expo Router \\
    TypeScript & TypeScript \\
    Convex & Convex \\
    Better Auth & Better Auth \\
    \bottomrule
  \end{tabular}
\end{center}

% =================================================
\section{Expo Router}
% =================================================

Expo Router administra la navegación de la aplicación móvil y pertenece a
Presentación.

Las pantallas estarán representadas mediante archivos dentro de:

\begin{lstlisting}
app/
\end{lstlisting}

Por ejemplo:

\begin{lstlisting}
app/
|
+-- _layout.tsx
+-- index.tsx
+-- login.tsx
|
+-- usuarios/
    |
    +-- index.tsx
    +-- [id].tsx
\end{lstlisting}

Expo Router se ocupa de:

\begin{itemize}
  \item navegación;
  \item layouts;
  \item parámetros;
  \item deep links;
  \item organización de pantallas.
\end{itemize}

La lógica de negocio no debe implementarse dentro del router. Las pantallas
delegan las operaciones en la Aplicación del cliente o en el cliente de
Infraestructura que invoca la API pública.

% =================================================
\section{Monorepo y un solo backend}
% =================================================

El backend Convex vive únicamente en \texttt{convex/}, en la raíz del monorepo.
TI4 \textbf{no debe crear otra carpeta Convex ni copiar funciones del backend}.

La arquitectura correcta es:

\begin{lstlisting}
Monorepo
|
 +-- convex/                         Backend
 |   +-- presentation/               Entrada pública Convex
 |   +-- application/                Casos de uso
 |   +-- domain/                     Reglas de negocio
 |   +-- infrastructure/             Adaptadores y persistencia
 |   +-- schema.ts
 |
 +-- apps/
     +-- web/                        Presentación Web
     +-- mobile/                     Presentación Mobile
\end{lstlisting}

Así existe una sola implementación del servidor y ambas aplicaciones consumen
la misma API pública. Cuando el producto crezca, organiza cada capa por
funcionalidad si ayuda a mantenerla clara; no mezcles responsabilidades para
evitar crear un archivo nuevo.

% =================================================
\section{Tipos compartidos dentro del monorepo}
% =================================================

Convex genera una representación TypeScript de su API pública. En el monorepo,
Web y Mobile acceden a esa representación desde la misma base de código.

El flujo de tipos es:

\begin{lstlisting}
convex/
     |
     | define funciones y genera tipos
     v
convex/_generated/api
     |                 |
     v                 v
apps/web          apps/mobile
\end{lstlisting}

Web y Mobile usan los tipos generados por Convex mediante un alias o un paquete
interno definido por el monorepo. La configuración del alias puede variar, pero
no dupliques la implementación del backend ni mantengas un contrato generado
por separado para cada aplicación.

Estos tipos pertenecen al borde entre Infraestructura y Presentación. No son el
modelo de Dominio ni deben hacer que las entidades del negocio dependan de
\texttt{convex/\_generated}. Cuando necesites contratos más estables, usa
\texttt{packages/shared/} y mantén separadas las reglas del negocio.

Declara la dependencia de Convex en los workspaces correspondientes e instala
el paquete desde la raíz:

\begin{lstlisting}
bun add convex --cwd apps/web
bun add convex --cwd apps/mobile
\end{lstlisting}

Ejecuta el backend desde la raíz:

\begin{lstlisting}
bunx convex dev
\end{lstlisting}

Mantén \texttt{convex/\_generated/} según la política del proyecto para archivos
generados. Si Convex puede regenerarla automáticamente en cada entorno, no
crees una segunda copia manual dentro de \texttt{apps/mobile/}.

\subsection{Regla}

Web y Mobile consumen la API pública definida en \texttt{convex/}. TI4 puede
consumir las funciones del backend y sus tipos, pero no modifica funciones
Convex fuera de su responsabilidad ni duplica en Mobile los casos de uso del
backend.

% =================================================
\section{Estructura del monorepo}
% =================================================

La estructura inicial puede ser:

\begin{lstlisting}
proyecto/
|
 +-- apps/
 |   |
 |   +-- web/
 |   |   +-- src/
 |   |       +-- presentation/
 |   |       +-- application/
 |   |       +-- infrastructure/
 |   |   +-- public/
 |   |   +-- package.json
 |   |   +-- vite.config.ts
 |   |
 |   +-- mobile/
 |       +-- app/
 |       +-- src/
 |           +-- application/
 |           +-- infrastructure/
 |       +-- package.json
 |       +-- app.json
|
 +-- packages/
 |   +-- ui/                 Componentes compartidos si corresponde
 |   +-- shared/             Tipos o utilidades compartidas si corresponde
|
 +-- convex/
 |   +-- _generated/
 |   +-- presentation/
 |   +-- application/
 |   +-- domain/
 |   +-- infrastructure/
 |   +-- betterAuth/
 |   +-- schema.ts
 |   +-- auth.config.ts
 |   +-- convex.config.ts
 |   +-- http.ts
 |   +-- ...
|
 +-- .github/
 |   +-- workflows/
 |
 +-- .env.example
 +-- .gitignore
 +-- package.json
 +-- bun.lock
 +-- tsconfig.json
 +-- vercel.json
 +-- README.md
\end{lstlisting}

Las carpetas de las capas marcan límites de responsabilidad; no obligan a crear
archivos vacíos. Una funcionalidad pequeña puede comenzar con pocos archivos si
cada uno tiene una responsabilidad clara y conserva la dirección de
dependencias.

Los directorios \texttt{packages/ui/} y \texttt{packages/shared/} son
opcionales. Créelos solo cuando exista una responsabilidad compartida real; no
los incorpores como capas vacías por anticipación.

El directorio del backend:

\begin{lstlisting}
convex/_generated/
\end{lstlisting}

Convex genera este directorio desde su configuración; no lo modifiques
manualmente.

% =================================================
La aplicación móvil usa el backend y los paquetes compartidos desde este
monorepo; no tiene una copia local del backend ni un archivo de contrato
externo.

% =================================================
\section{Setup inicial del monorepo}
% =================================================

Ejecuta todos los comandos siguientes desde la raíz del monorepo.

\subsection{1. Crear la raíz y los workspaces}

Declara en la raíz los workspaces de las aplicaciones y de los paquetes
internos:

\begin{lstlisting}
{
  "name": "proyecto",
  "private": true,
  "workspaces": [
    "apps/*",
    "packages/*"
  ]
}
\end{lstlisting}

\subsection{2. Crear aplicación Web}

\begin{lstlisting}
bun create vite@latest apps/web -- --template react-ts
bun install
\end{lstlisting}

\subsection{3. Crear aplicación Mobile}

\begin{lstlisting}
bunx create-expo-app@latest apps/mobile
bun install
\end{lstlisting}

La plantilla estándar de Expo ya incorpora Expo Router.

\subsection{4. Instalar Convex en la raíz}

Configura el backend Convex en la raíz del monorepo:

\begin{lstlisting}
bun add convex
bunx convex dev
\end{lstlisting}

La primera ejecución vincula o crea el proyecto Convex y genera:

\begin{lstlisting}
convex/
\end{lstlisting}

\subsection{5. Crear los límites de las capas}

Después de crear el backend, establece los límites de las capas antes de
agregar lógica de negocio:

\begin{lstlisting}
convex/
+-- presentation/       Funciones públicas de Convex
+-- application/        Casos de uso
+-- domain/              Entidades y reglas de negocio
+-- infrastructure/      Convex, autenticación y servicios externos

apps/web/src/
+-- presentation/
+-- application/
+-- infrastructure/

apps/mobile/
+-- app/                 Rutas y pantallas
+-- src/application/
+-- src/infrastructure/
\end{lstlisting}

No crees archivos vacíos en cada carpeta. Decide primero dónde vive la
responsabilidad y luego incorpora el código mínimo de la funcionalidad.

\subsection{6. Instalar TanStack Router en Web}

\begin{lstlisting}
bun add @tanstack/react-router --cwd apps/web
bun add -d @tanstack/router-plugin --cwd apps/web
bun add -d @tanstack/react-router-devtools --cwd apps/web
\end{lstlisting}

Agrega en Vite el plugin de TanStack Router antes del plugin de React.

Ejemplo:

\begin{lstlisting}
import { defineConfig } from 'vite'
import react from '@vitejs/plugin-react'
import { tanstackRouter } from '@tanstack/router-plugin/vite'

export default defineConfig({
  plugins: [
    tanstackRouter({
      target: 'react',
      autoCodeSplitting: true,
    }),
    react(),
  ],
})
\end{lstlisting}

\subsection{7. Instalar Better Auth en el backend}

\begin{lstlisting}
bun add better-auth @convex-dev/better-auth
\end{lstlisting}

Realiza la integración dentro del backend Convex y sigue la estructura oficial
de Better Auth para Convex.

Entre otros, la integración usa:

\begin{lstlisting}
convex/
|
+-- auth.config.ts
+-- convex.config.ts
+-- http.ts
|
+-- betterAuth/
    |
    +-- convex.config.ts
    +-- auth.ts
    +-- schema.ts
    +-- adapter.ts
\end{lstlisting}

Genera el esquema de Better Auth con sus herramientas; evita mantenerlo
manualmente cuando sea posible.

% =================================================
\section{Dependencias de la aplicación Mobile}
% =================================================

TI4 usa el cliente Convex y los tipos generados desde la carpeta compartida del
backend. No crea ni despliega una segunda implementación.

\subsection{1. Instalar Better Auth para Expo}

\begin{lstlisting}
bun add better-auth @better-auth/expo --cwd apps/mobile
bun add expo-network expo-secure-store --cwd apps/mobile
\end{lstlisting}

Better Auth usa Secure Store para almacenar la sesión de forma apropiada en la
aplicación móvil.

Si el proyecto usa proveedores sociales, también puede necesitar dependencias
de Expo como linking o web browser.

% =================================================
\section{Conexión de la aplicación Web a Convex}
% =================================================

La aplicación web usa una instancia de:

\begin{lstlisting}
ConvexReactClient
\end{lstlisting}

configurada con la URL del deployment.

La variable pública de Vite es:

\begin{lstlisting}
VITE_CONVEX_URL=
\end{lstlisting}

Dentro del frontend:

\begin{lstlisting}
const convex = new ConvexReactClient(
  import.meta.env.VITE_CONVEX_URL
)
\end{lstlisting}

Provee el cliente en la raíz de la aplicación.

% =================================================
\section{Conexión de la aplicación Mobile a Convex}
% =================================================

Expo usa el mismo cliente React de Convex.

La variable de Expo es:

\begin{lstlisting}
EXPO_PUBLIC_CONVEX_URL=
\end{lstlisting}

Configura en la raíz de Expo Router:

\begin{lstlisting}
const convex = new ConvexReactClient(
  process.env.EXPO_PUBLIC_CONVEX_URL!
)
\end{lstlisting}

Luego, provee el cliente a la aplicación mediante el provider correspondiente.

Ambas aplicaciones usan el mismo deployment:

\begin{lstlisting}
Web
  |
  +----\
        \
         > mismo deployment Convex
        /
Mobile /
\end{lstlisting}

% =================================================
\section{URLs de Convex}
% =================================================

Convex usa URLs con responsabilidades distintas.

El proyecto distingue:

\begin{description}
  \item[Convex URL:] conexión del cliente React con las funciones Convex.
  \item[Convex Site URL:] endpoints HTTP, incluyendo rutas utilizadas por
    Better Auth cuando corresponda.
\end{description}

Los clientes pueden necesitar variables equivalentes a:

\begin{lstlisting}
VITE_CONVEX_URL=
VITE_CONVEX_SITE_URL=
\end{lstlisting}

para Web, y:

\begin{lstlisting}
EXPO_PUBLIC_CONVEX_URL=
EXPO_PUBLIC_CONVEX_SITE_URL=
\end{lstlisting}

para mobile.

Estas URLs localizan servicios públicos.

\textbf{No son el lugar donde deben almacenarse secretos.}

% =================================================
\section{Variables de entorno}
% =================================================

Distingue las variables públicas de los secretos.

\subsection{Frontend Web}

Vite expone al código cliente variables con prefijo:

\begin{lstlisting}
VITE_
\end{lstlisting}

Por ejemplo:

\begin{lstlisting}
VITE_CONVEX_URL=
\end{lstlisting}

Toda variable \texttt{VITE\_} es visible para el usuario final.

\subsection{Aplicación móvil}

Expo expone al código cliente las variables con prefijo:

\begin{lstlisting}
EXPO_PUBLIC_
\end{lstlisting}

Por ejemplo:

\begin{lstlisting}
EXPO_PUBLIC_CONVEX_URL=
\end{lstlisting}

Toda variable \texttt{EXPO\_PUBLIC\_} es pública.

\subsection{Backend}

Configura en el entorno de Convex los secretos que usan sus funciones.

Por ejemplo:

\begin{lstlisting}
BETTER_AUTH_SECRET
\end{lstlisting}

Estos valores no deben estar:

\begin{itemize}
  \item dentro del código;
  \item dentro del frontend;
  \item dentro de variables públicas;
  \item versionados en Git.
\end{itemize}

% =================================================
\section{Archivos de entorno}
% =================================================

No subas al monorepo archivos locales con valores específicos de un
desarrollador.

Cada aplicación puede mantener su archivo local dentro de su directorio:

\begin{lstlisting}
apps/web/.env.local
apps/mobile/.env.local
\end{lstlisting}

Mantén en el monorepo ejemplos de configuración sin secretos:

\begin{lstlisting}
apps/web/.env.example
apps/mobile/.env.example
\end{lstlisting}

Incluye las variables necesarias, pero ningún secreto.

Ejemplo TI2:

\begin{lstlisting}
VITE_CONVEX_URL=
VITE_CONVEX_SITE_URL=
\end{lstlisting}

Ejemplo TI4:

\begin{lstlisting}
EXPO_PUBLIC_CONVEX_URL=
EXPO_PUBLIC_CONVEX_SITE_URL=
\end{lstlisting}

% =================================================
\section{Ejecución local del monorepo}
% =================================================

Durante el desarrollo ejecuta al menos dos procesos desde la raíz del monorepo.

\subsection{Terminal 1: Convex}

\begin{lstlisting}
bunx convex dev
\end{lstlisting}

Este proceso sincroniza los cambios del backend con el deployment de desarrollo.

\subsection{Terminal 2: Web}

\begin{lstlisting}
bun --cwd apps/web run dev
\end{lstlisting}

Los procesos se relacionan así:

\begin{lstlisting}
Terminal 1                 Terminal 2

bunx convex dev            bun --cwd apps/web run dev
      |                          |
      v                          v
Backend Convex             React + Vite
\end{lstlisting}

% =================================================
\section{Ejecución local de Mobile}
% =================================================

Ejecuta la aplicación móvil mediante Expo.

\begin{lstlisting}
bun --cwd apps/mobile run start
\end{lstlisting}

o equivalentemente:

\begin{lstlisting}
cd apps/mobile && bunx expo start
\end{lstlisting}

Desde esta sesión puedes ejecutar la aplicación en:

\begin{itemize}
  \item Android;
  \item iOS cuando el entorno lo permita;
  \item dispositivo físico;
  \item emulador o simulador.
\end{itemize}

La aplicación Mobile debe conectarse al deployment de desarrollo de Convex
definido para el proyecto.

No necesitas ejecutar el backend para modificar exclusivamente Mobile, pero sí
debes usar un deployment válido para probar la integración.

% =================================================
\section{Ambientes}
% =================================================

Distingue al menos:

\begin{description}
  \item[Development:] programación y pruebas durante el desarrollo.
  \item[Production:] versión estable utilizada para las demostraciones y
    entregas.
\end{description}

En Mobile también puede existir:

\begin{description}
  \item[Preview:] build destinado a pruebas del equipo.
\end{description}

Usa configuraciones separadas cuando sea necesario.

% =================================================
\section{Build de la aplicación Web}
% =================================================

Antes de integrar una versión, genera su build de producción.

\begin{lstlisting}
bun --cwd apps/web run build
\end{lstlisting}

Una aplicación que solo funciona con:

\begin{lstlisting}
bun --cwd apps/web run dev
\end{lstlisting}

pero cuyo build falla no está lista para integrar.

% =================================================
\section{Despliegue del backend Convex}
% =================================================

Durante el desarrollo usa:

\begin{lstlisting}
bunx convex dev
\end{lstlisting}

Para desplegar funciones, esquema e índices en producción usa:

\begin{lstlisting}
bunx convex deploy
\end{lstlisting}

TI2 es responsable del deployment de producción y lo ejecuta desde la raíz del
monorepo.

No todos los integrantes deben manipular producción manualmente.

% =================================================
\section{Despliegue de la aplicación Web}
% =================================================

Vercel despliega la aplicación web.

Conecta Vercel al único repositorio GitHub y configura \texttt{apps/web/} como
directorio de la aplicación. Las órdenes de instalación y build deben respetar
el \texttt{bun.lock} de la raíz.

El flujo general es:

\begin{lstlisting}
GitHub
    |
    v
   main
    |
    v
 Vercel
    |
    +---- Build apps/web
    |
    +---- Deploy Convex
    |
    v
Producción
\end{lstlisting}

Para integrar el despliegue de Convex con Vercel, configura un deploy key de
producción como variable segura:

\begin{lstlisting}
CONVEX_DEPLOY_KEY
\end{lstlisting}

El build puede ejecutar:

\begin{lstlisting}
bunx convex deploy --cmd 'bun --cwd apps/web run build'
\end{lstlisting}

Así reduces el riesgo de publicar un frontend incompatible con las funciones
Convex correspondientes.

% =================================================
\section{Configuración SPA en Vercel}
% =================================================

TanStack Router administra las rutas en el cliente.

Por eso, al solicitar directamente una ruta como:

\begin{lstlisting}
/usuarios/123
\end{lstlisting}

el hosting debe entregar la aplicación y permitir que TanStack Router resuelva
la ruta.

El repositorio puede incluir:

\begin{lstlisting}
vercel.json
\end{lstlisting}

con una regla equivalente a:

\begin{lstlisting}
{
  "rewrites": [
    {
      "source": "/(.*)",
      "destination": "/index.html"
    }
  ]
}
\end{lstlisting}

% =================================================
\section{Build de la aplicación móvil}
% =================================================

EAS Build genera los builds para instalar. Configura EAS con
\texttt{apps/mobile/}
como directorio de la aplicación y comparte el lockfile de la raíz cuando CI se
ejecute desde el monorepo.

Primero configura EAS para el proyecto.

Una configuración simplificada de:

\begin{lstlisting}
eas.json
\end{lstlisting}

puede incluir:

\begin{lstlisting}
{
  "build": {
    "development": {
      "developmentClient": true
    },
    "preview": {
      "distribution": "internal"
    },
    "production": {}
  }
}
\end{lstlisting}

% =================================================
\section{Build Preview de Mobile}
% =================================================

Durante el curso no es necesario publicar cada versión en Google Play o App
Store.

Para las pruebas usa preferentemente un build interno.

Android:

\begin{lstlisting}
eas build --platform android --profile preview
\end{lstlisting}

iOS:

\begin{lstlisting}
eas build --platform ios --profile preview
\end{lstlisting}

EAS genera un build que puedes compartir con integrantes y personas autorizadas
para pruebas.

En Android, este mecanismo permite instalar un APK sin publicarlo en Play Store.

% =================================================
\section{Build de producción Mobile}
% =================================================

Cuando necesites generar una versión de producción:

\begin{lstlisting}
eas build --platform android
\end{lstlisting}

o:

\begin{lstlisting}
eas build --platform ios
\end{lstlisting}

También puedes usar:

\begin{lstlisting}
eas build --platform all
\end{lstlisting}

La publicación efectiva en tiendas es independiente del proceso de desarrollo.

% =================================================
\section{Flujos de despliegue}
% =================================================

\subsection{TI2}

\begin{lstlisting}
Monorepo GitHub
  |
  v
main
  |
  v
Vercel
  |
  +--> React/Vite
  |
  +--> Convex Production
  |
  v
Web desplegada
\end{lstlisting}

\subsection{TI4}

\begin{lstlisting}
Monorepo GitHub
  |
  v
main
  |
  v
EAS Build
  |
  +--> Preview
  |
  +--> Production
  |
  v
Aplicación para instalar
\end{lstlisting}

% =================================================
\section{Integración Continua}
% =================================================

GitHub Actions verifica los cambios del monorepo antes de integrarlos.

Instala desde la raíz para utilizar el lockfile único:

\begin{lstlisting}
bun install --frozen-lockfile
\end{lstlisting}

Cuando existan los scripts correspondientes, ejecuta como mínimo estas
verificaciones sobre los componentes afectados:

\begin{enumerate}
  \item instalación de dependencias desde la raíz;
  \item chequeo de TypeScript;
  \item lint;
  \item pruebas unitarias del Dominio y de los casos de uso;
  \item pruebas de integración de los adaptadores de Infraestructura cuando
    corresponda;
  \item pruebas de Presentación disponibles;
  \item build de Web, Mobile o backend según las rutas modificadas.
\end{enumerate}

La secuencia de CI es:

\begin{lstlisting}
Pull Request
     |
     v
bun install --frozen-lockfile
     |
     v
Type Check
     |
     v
Lint
     |
     v
Tests
     |
     v
Build
     |
     v
Code Review
     |
     v
Merge
\end{lstlisting}

El despliegue no reemplaza CI.

% =================================================
\section{Responsabilidades de TI2}
% =================================================

TI2 es responsable técnicamente de:

\begin{itemize}
  \item desarrollar la aplicación web;
  \item mantener React y TanStack Router;
  \item implementar la Presentación pública del backend mediante funciones
    Convex delgadas;
  \item implementar los casos de uso de Aplicación;
  \item definir y mantener las entidades y reglas del Dominio;
  \item implementar los adaptadores de Infraestructura;
  \item definir y mantener el modelo de datos;
  \item implementar Better Auth;
  \item implementar reglas de autorización;
  \item desplegar el backend;
  \item desplegar la aplicación web;
  \item mantener disponibles los servicios requeridos por TI4;
  \item mantener estable la API pública consumida por Mobile;
  \item informar cambios incompatibles;
  \item mantener actualizados los tipos generados y los paquetes compartidos
    cuando corresponda.
\end{itemize}

TI2 es propietario técnico del backend compartido.

% =================================================
\section{Responsabilidades de TI4}
% =================================================

TI4 es responsable técnicamente de:

\begin{itemize}
  \item desarrollar la aplicación móvil mediante Expo;
  \item organizar la navegación mediante Expo Router;
  \item implementar la Presentación Mobile;
  \item organizar la Aplicación específica de la interfaz cuando una pantalla
    requiera más de una llamada al backend;
  \item integrar la Infraestructura cliente con Convex;
  \item integrar los flujos de Better Auth;
  \item consumir los tipos generados desde la configuración compartida;
  \item detectar y reportar problemas de integración;
  \item generar builds mediante EAS;
  \item evitar duplicar lógica perteneciente al backend.
\end{itemize}

TI4 también mantiene las responsabilidades de gestión, coordinación,
seguimiento y control de calidad definidas en el documento metodológico general.

% =================================================
\section{Frontera entre TI2 y TI4}
% =================================================

La frontera técnica entre ambos equipos está formada por la API pública de
Presentación del backend, los casos de uso que expone, el sistema de
autenticación compartido y los paquetes internos que consumen ambas
aplicaciones.

TI4 no debe:

\begin{itemize}
  \item crear una segunda base de datos;
  \item crear una segunda implementación del backend;
  \item copiar las funciones Convex dentro de \texttt{apps/mobile/};
  \item modificar directamente producción de TI2;
  \item duplicar reglas de negocio únicamente para hacer funcionar Mobile.
\end{itemize}

TI2 debe documentar suficientemente las funciones públicas que consume Mobile.

Como mínimo, documenta:

\begin{itemize}
  \item propósito;
  \item argumentos;
  \item resultado;
  \item requerimientos de autenticación;
  \item errores relevantes.
\end{itemize}

% =================================================
\section{Cambios que afectan a ambos equipos}
% =================================================

Coordina mediante Linear cualquier cambio que modifique:

\begin{itemize}
  \item nombre de una función pública;
  \item argumentos de una función;
  \item tipo de retorno;
  \item comportamiento de un caso de uso;
  \item reglas de autenticación;
  \item reglas de autorización;
  \item regla o política del Dominio;
  \item estructura de datos utilizada por ambos clientes;
  \item comportamiento funcional esperado por la aplicación Mobile;
  \item un paquete compartido utilizado por Web y Mobile.
\end{itemize}

Registra estas dependencias en Linear.

Por ejemplo:

\begin{lstlisting}
TI2-42 Cambiar flujo de autenticación
              |
              +-- bloquea -->
                              TI4-18
                              Adaptar login mobile
\end{lstlisting}

% =================================================
\section{Qué corresponde a cada capa}
% =================================================

\subsection*{Presentación}

React, React Native, TanStack Router, Expo Router y las funciones públicas de
Convex forman los puntos de entrada de esta capa. Presentación se ocupa de
recibir eventos, validar datos de entrada y mostrar estados de carga, éxito y
error.

No contiene acceso directo a la base de datos ni reglas de negocio compartidas.

\subsection*{Aplicación}

Los casos de uso representan acciones con significado para el producto, por
ejemplo \textit{Obtener perfil}, \textit{Registrar usuario} o
\textit{Reservar recurso}. Esta capa coordina autenticación, autorización,
validaciones de negocio y persistencia mediante puertos.

Un caso de uso no debe importar React, Expo, Convex Client ni componentes de
interfaz.

\subsection*{Dominio}

El Dominio contiene entidades, objetos de valor, invariantes y políticas del
negocio. Es la única capa que debe ser la fuente de verdad para las reglas que
deben cumplirse independientemente de Web o Mobile.

Es código TypeScript independiente de frameworks. No recibe directamente
documentos de Convex ni tipos de navegación.

\subsection*{Infraestructura}

Convex, el esquema, los repositorios, Better Auth, Convex Client y las
integraciones con servicios externos forman adaptadores de Infraestructura.
Esta capa conoce los detalles técnicos y traduce esos detalles a los puertos
que consumen Aplicación y Dominio.

\subsection*{Herramientas de soporte}

Vite, Expo, EAS Build, Vercel, Bun y GitHub Actions son herramientas de
desarrollo, construcción o despliegue. No son capas de la arquitectura y no
deben utilizarse como ubicación para reglas del negocio.

% =================================================
\section{Ejemplo de distribución de responsabilidades}
% =================================================

Supón que debes implementar:

\begin{center}
  \textit{``Mostrar el perfil de un usuario autenticado.''}
\end{center}

La responsabilidad se distribuye así:

\begin{lstlisting}
Presentación
    |
    +--> Web o Mobile solicita el perfil
    |
    v
Aplicación
    |
    +--> caso de uso Obtener perfil
    +--> solicita identidad y repositorio
    |
    v
Dominio
    |
    +--> valida la política de acceso
    |
    v
Infraestructura
    |
    +--> Better Auth identifica al usuario
    +--> Repositorio Convex obtiene el perfil
    |
    +-------------+
    |             |
    v             v
   Web          Mobile
  React       React Native
\end{lstlisting}

El frontend decide \textbf{cómo mostrar} la información. Aplicación y Dominio
deciden \textbf{qué información puede obtener} el usuario. Infraestructura
implementa la forma concreta de obtenerla.

% =================================================
\section{Proceso para una funcionalidad compartida}
% =================================================

Una funcionalidad que afecte a ambos equipos sigue este flujo:

\begin{enumerate}
  \item TI4 hace el triage y los equipos identifican las tareas necesarias.
  \item TI2 define el caso de uso y las reglas del Dominio.
  \item TI2 implementa los puertos y adaptadores de Infraestructura necesarios.
  \item TI2 expone una función pública de Convex como entrada del caso de uso.
  \item TI2 y TI4 integran el cambio mediante uno o más Pull Requests
    relacionados.
  \item Regeneran o actualizan los tipos públicos de Convex.
  \item TI4 implementa el consumo desde la aplicación Mobile.
  \item Ambos equipos prueban la integración.
  \item Los cambios llegan a los procesos de despliegue de Web y Mobile.
\end{enumerate}

La dependencia entre equipos queda así:

\begin{lstlisting}
Linear
   |
   +-------------------+
   |                   |
   v                   v
Issue TI2          Issue TI4
   |                   |
   v                   |
Caso de uso             |
   |                   |
   v                   |
API publica -------+    |
   |               |    |
   v               |    |
Tipos generados ---+----+
                       |
                       v
                  apps/mobile
\end{lstlisting}

% =================================================
\section{Principios técnicos}
% =================================================

Durante el desarrollo rigen estos principios:

\begin{enumerate}
  \item existe un solo backend;
  \item existe una sola fuente de verdad para los datos;
  \item Web y Mobile consumen los mismos servicios;
  \item las dependencias respetan la dirección Presentación $\rightarrow$
    Aplicación $\rightarrow$ Dominio;
  \item el Dominio no depende de frameworks ni proveedores concretos;
  \item los casos de uso coordinan las operaciones y las funciones públicas
    de Convex se mantienen delgadas;
  \item la Infraestructura implementa puertos sin trasladar reglas de negocio
    a los repositorios;
  \item los secretos nunca se almacenan en código cliente;
  \item las reglas de autorización se validan en el backend;
  \item la API pública y los paquetes compartidos deben ser explícitos;
  \item los cambios incompatibles deben coordinarse antes de integrarse;
  \item el código generado necesario para compilar debe mantenerse
    correctamente versionado;
  \item una aplicación que no puede generar su build no está terminada;
  \item los entornos de desarrollo y producción deben mantenerse separados.
\end{enumerate}

% =================================================
\section{Resumen de arquitectura}
% =================================================

El stack tecnológico puede resumirse de la siguiente manera:

\begin{lstlisting}
       Web / Mobile
       Presentación
             |
             v
       Aplicación
       Casos de uso
             |
             v
         Dominio
     Reglas de negocio
             ^
             |
      Infraestructura
   Convex + Better Auth
\end{lstlisting}

React, React Native, TanStack Router y Expo Router se ubican principalmente en
Presentación. TypeScript aparece en todas las capas. Vite, Expo, EAS Build,
Vercel y GitHub Actions soportan el desarrollo y despliegue, pero no contienen
reglas de negocio.

El flujo de desarrollo y despliegue es:

\begin{lstlisting}
                       Linear
                         |
                         v
                      Issue
                         |
                         v
                  Feature Branch
                         |
                         v
                   Pull Request
                         |
                 +-------+-------+
                 |               |
                 v               v
                CI          Code Review
                 |               |
                 +-------+-------+
                         |
                         v
                        main
                         |
             +-----------+-----------+
             |                       |
             v                       v
         TI2 Deploy              TI4 Build
             |                       |
       Vercel + Convex            EAS Build
             |                       |
             v                       v
            Web                  Mobile App
\end{lstlisting}

% =================================================
\section{Referencias oficiales}
% =================================================

Las configuraciones específicas pueden cambiar con nuevas versiones de las
herramientas. Antes de cambiar una configuración importante, consulta la
documentación oficial correspondiente.

\begin{itemize}

  \item TanStack Router:
    \url{https://tanstack.com/router/latest}

  \item TanStack Router con Vite:
    \href{https://tanstack.com/router/latest/docs/installation/with-vite}{Documentación
    de TanStack Router con Vite}

  \item Convex:
    \url{https://docs.convex.dev/}

  \item Convex con React:
    \url{https://docs.convex.dev/client/react/overview}

  \item Convex con React Native:
    \url{https://docs.convex.dev/quickstart/react-native}

  \item Convex en Vercel:
    \url{https://docs.convex.dev/production/hosting/vercel}

  \item Better Auth:
    \url{https://better-auth.com/}

  \item Better Auth con Convex:
    \url{https://better-auth.com/docs/integrations/convex}

  \item Better Auth con Expo:
    \url{https://better-auth.com/docs/integrations/expo}

  \item Expo:
    \url{https://docs.expo.dev/}

  \item EAS Build:
    \url{https://docs.expo.dev/build/introduction/}

  \item EAS Internal Distribution:
    \url{https://docs.expo.dev/build/internal-distribution/}

\end{itemize}

\end{document}
